\chapter{Manual de usuario}

\section{Uso general}

El usuario inicia una ejecución proporcionando una descripción del problema de
toma de decisiones. El sistema genera investigación, formulaciones,
especificaciones y modelos en etapas sucesivas.

\section{Revisiones}

En los puntos de revisión, el usuario debe aprobar o rechazar artefactos. La
revisión es importante porque determina qué conocimiento puede pasar a memoria.

\section{Modelo generado}

Un modelo aceptado puede ser usado por el laboratorio virtual si implementa:

\begin{lstlisting}[language=Python]
def decide(self, perception: dict) -> Action
def update(self, action, reward, new_perception) -> None
def get_state(self) -> dict
\end{lstlisting}

\section{Errores}

Si falla un servicio opcional de memoria, la ejecución puede continuar en modo
degradado. Si fallan PostgreSQL o MinIO, la ejecución no puede completarse
correctamente porque son necesarios para persistencia y artefactos.
